%==============================================================================
\section{Act IV: The Photo (\texttt{Frame.h/.cc})}
%==============================================================================

The \texttt{Frame} class represents a single image capture. It is the fundamental data unit that flows through the system.

\subsection{What Does a Frame Contain?}

\begin{lstlisting}[caption={Key Frame members}]
class Frame {
    // Identity
    long unsigned int mnId;       // Unique frame ID
    double mTimeStamp;            // When captured
    
    // Features (the interesting points)
    int N;                        // Number of keypoints
    vector<cv::KeyPoint> mvKeys;  // 2D keypoint locations
    vector<cv::KeyPoint> mvKeysUn;  // Undistorted keypoints
    cv::Mat mDescriptors;         // Feature descriptors (Nx32 for ORB)
    
    // 3D associations
    vector<MapPoint*> mvpMapPoints;  // Pointer to 3D point for each keypoint
    vector<bool> mvbOutlier;      // Outlier flags
    
    // Camera
    cv::Mat mK;                   // Calibration matrix
    cv::Mat mDistCoef;            // Distortion coefficients
    GeometricCamera* mpCamera;    // Camera model
    
    // Pose
    Sophus::SE3f mTcw;            // Camera pose (world to camera)
    bool mbHasPose;               // Is pose valid?
    
    // Optimization structures
    vector<size_t> mGrid[64][48]; // Spatial grid for fast search
    DBoW2::BowVector mBowVec;     // Bag-of-Words representation
};
\end{lstlisting}

\subsection{The Frame Constructor Journey}

When a new Frame is created, this happens:

\begin{enumerate}
    \item \textbf{Assign ID}: \texttt{mnId = nNextId++}
    \item \textbf{Extract Features}: Call ORBextractor (or XFeat)
    \item \textbf{Undistort Keypoints}: Correct for lens distortion
    \item \textbf{Initialize Associations}: \texttt{mvpMapPoints} all set to NULL
    \item \textbf{Assign to Grid}: Organize keypoints into spatial cells
    \item \textbf{Setup Scale Pyramid}: Prepare multi-scale information
\end{enumerate}

\subsection{The Different Constructors}

Frame has multiple constructors for different scenarios:

\begin{table}[H]
\centering
\begin{tabular}{@{}ll@{}}
\toprule
\textbf{Constructor} & \textbf{Use Case} \\
\midrule
\texttt{Frame()} & Default, empty frame \\
\texttt{Frame(const Frame\&)} & Copy constructor \\
\texttt{Frame(imLeft, imRight, ...)} & Stereo cameras \\
\texttt{Frame(imGray, imDepth, ...)} & RGB-D cameras \\
\texttt{Frame(imGray, timestamp, ...)} & Monocular camera \\
\texttt{Frame(imGray, keypoints, descriptors, ...)} & Pure XFeat mode \\
\bottomrule
\end{tabular}
\caption{Frame constructor variants}
\end{table}

\subsection{Monocular vs. XFeat Constructor}

The standard constructor calls \texttt{ExtractORB()}, which is slow.

Your custom \textbf{Pure XFeat constructor} skips this:

\begin{table}[H]
\centering
\begin{tabular}{@{}lcc@{}}
\toprule
\textbf{Step} & \textbf{Standard ORB} & \textbf{Pure XFeat} \\
\midrule
Feature Extraction & Calls \texttt{ExtractORB()} & Skipped (pre-extracted) \\
Undistortion & Calls \texttt{UndistortKeyPoints()} & Skipped (pre-undistorted) \\
Descriptors & 256-bit binary & 64-float vectors \\
Scale Levels & From ORBextractor (8 levels) & Manually set (3 levels) \\
\bottomrule
\end{tabular}
\caption{Comparison of Frame constructors}
\end{table}

\subsection{The Spatial Grid}

To speed up feature matching, the Frame divides the image into a grid:

\begin{lstlisting}[caption={Spatial grid for fast lookup}]
#define FRAME_GRID_COLS 64
#define FRAME_GRID_ROWS 48

// Each cell contains indices of keypoints in that region
vector<size_t> mGrid[FRAME_GRID_COLS][FRAME_GRID_ROWS];
\end{lstlisting}

This allows fast queries like: ``Find all features within 20 pixels of point (x, y)'' without scanning all N keypoints.

\subsection{Key Methods}

\begin{lstlisting}[caption={Important Frame methods}]
// Set/Get the camera pose
void SetPose(const Sophus::SE3f &Tcw);
Sophus::SE3f GetPose() const;

// Compute Bag-of-Words (for loop closure)
void ComputeBoW();

// Find features in a region (uses the grid)
vector<size_t> GetFeaturesInArea(float x, float y, float r);

// Check if a MapPoint is visible
bool isInFrustum(MapPoint* pMP, float viewingCosLimit);
\end{lstlisting}

\subsection{Frame vs KeyFrame}

\begin{table}[H]
\centering
\begin{tabular}{@{}lcc@{}}
\toprule
\textbf{Aspect} & \textbf{Frame} & \textbf{KeyFrame} \\
\midrule
Lifespan & Temporary & Permanent \\
Storage & Only current/last & All in map \\
Connections & None & Covisibility graph \\
Optimization & Pose only & Pose + structure \\
Loop closure & Cannot trigger & Can trigger \\
Count per sequence & Thousands & Hundreds \\
\bottomrule
\end{tabular}
\caption{Frames are temporary; KeyFrames are permanent}
\end{table}
